\chapter*{Preface}
\addcontentsline{toc}{chapter}{Preface}
\markboth{Preface}{Preface}

Statistical physics learned to solve models on networks by pretending the
networks were trees. That assumption is what makes the method work. If a network has no short loops, then cutting one node out of it
leaves the pieces that hang off that node independent of one another, and a
problem defined over the whole network --- every spin coupled to every
neighbour, every constraint tangled with every other --- collapses into a single
equation for what one branch tells the node above it. That equation is the
cavity recursion, or belief propagation, or the Bethe--Peierls approximation,
depending on which decade and which field one learned it in. From it come
critical temperatures, percolation thresholds, the size of a minimal vertex
cover, the number of solutions of a random satisfiability problem. It is one of
the most productive ideas in the statistical mechanics of disordered systems.

Real networks have short loops: the neighbours of a node are neighbours of one
another. Proteins act in complexes of four and five, not in pairs. A paper has six authors, and all
fifteen of the pairs among them are ``co-authorship'' links that the paper
created in one act. A chemical reaction consumes three species and produces
two. In each case the pairwise picture --- the graph, dots joined by lines ---
records the shadow of an interaction rather than the interaction itself, and it
records it in the form of a small dense knot of links that is exactly what the
tree assumption forbids. So the field has had, for twenty-five years, a method
that works and a systematic reason why it should not apply to the objects it is
applied to.

This book is about a way out that is simple once stated. If the trouble is a dense knot, stop treating it as many links and
start treating it as one thing. Call that thing a \emph{complex}, let a complex
contain nodes, and then --- this is the step that does the work --- let a
complex contain other complexes, to any depth one likes. The resulting object is
a \emph{complex hypergraph}, or chygraph: a set of complexes whose vertices are
themselves complexes. The definition is plain, and it covers a surprising range
of structures. A graph is a chygraph whose complexes all have two vertices.
A hypergraph is a chygraph with one layer of larger complexes. A multiplex
network, an interacting collection of hypergraphs, a graph decorated with
triangles and squares: all of them are chygraphs, differing only in what one
writes in the layer data. And a clustered graph becomes a chygraph by promoting
its dense motifs to complexes --- after which the structure \emph{over}
complexes can be treelike even though the graph underneath is not.

That is the bargain the book explores. Give up the graph as the fundamental
object, keep the tree, and the whole apparatus comes back. What comes back with
it is more than a rederivation. Percolation on every one of those structures
reduces to a single self-consistency map, and its threshold to the spectrum of
one tensor built from first moments. Push the same map one order further and the
critical amplitude follows in closed form. Replace the scalar message with a
measure and the Ising model follows; with an integer, the hitting set problem;
with three states, core percolation and the leaf-removal algorithm behind minimal
vertex covers; with a distribution over $q$ colours, graph colouring; with a
warning, satisfiability. And once a complex is allowed to hold other complexes
rather than only nodes, one can write down constraints on constraints, which no
flat encoding can hold without damage. Through all of it the recursion does not
change. Only what the
message carries changes, and what happens inside one complex, where the
Hamiltonian is summed exactly over an interior that may be as large and as
strongly coupled as one pleases. One $L\times L$ determinant, with $L$ the number
of layers and no dependence on the size of the complexes at all, serves every
model in this book.

Several of the results are counterintuitive. A triangle transmits more
magnetisation per neighbour than an edge does, and yet clustering \emph{lowers} the critical
temperature of the Ising model at fixed degree distribution --- because the
triangle also removes a branch, and the second effect is the larger. The
zero-temperature limit that everyone takes for hard-core covering problems is
exact on graphs and wrong on hypergraphs, giving a cover of one half where the
answer is one third, on an example so simple it can be checked by hand. A
complex of three or more nodes destroys the core-free phase of leaf removal
entirely, so that above cardinality two no minimal cover is ever certified. And
hyperbolic random graphs, whose clustering has defeated degree-based mean-field
theory for a decade, are predicted correctly once one stops applying the theory
to the graph and applies it to the chygraph of its cliques.

Where the bargain fails, the failure is structural rather than technical, so it
cannot be patched by working harder. A
chygraph is treelike only when its complexes meet in at most one atom. Cliques
in a geometric graph share edges; promoting them to complexes removes the loops
\emph{inside} motifs and exposes the loops \emph{between} them. Clustering is not one obstruction
but a hierarchy of them. Part~IV is the accounting: one chapter measures the
price, one runs the repair the literature already has and reports what it does
and does not recover, one merges the offending complexes and finds the condition
under which that is exact, and all three state plainly what is not yet done at
the level of the ensemble.

Who is this for? Anyone who has met the cavity method and wondered what to do
about clustering, and anyone who works with data whose interactions are
genuinely group-level --- proteins in complexes, papers with many authors,
reactions with many substrates, projects with many dependencies --- and has been
told to make a graph of it. The technical prerequisites are modest and are built
up in the first chapter: what percolation is, what a generating function does,
what the Ising model and vertex cover and the hitting set problem ask, and what
mean-field, cavity and replica theory mean. Readers who already know all of that
can start at Chapter~\ref{ch:chygraphs}. The longer calculations run in the text
where they are needed rather than in boxes set aside to be skipped, and each one
ends by naming the routine that computes it, so a reader who wants to run a
result rather than read it knows where to go. The book has a great many
pictures: nearly every idea here is about which things are grouped with which,
and grouping is easier to see in a drawing than in an equation.

A word on what is old and what is not, since the two are mixed.
Two papers stand behind parts of this book \citep{vazquez2023pre,vazquez2024comnet}:
the definition of a chygraph in Chapter~\ref{ch:chygraphs}, the percolation map
and threshold tensor of Chapter~\ref{ch:percolation}, and the interacting-network
constructions Chapter~\ref{ch:epidemics} builds on. Those chapters are expositions,
and they say so and cite as they go.

Everything else appears here for the first time. That includes the
clique-promotion measurements of Chapter~\ref{ch:data}; the critical
amplitude and the moment hierarchy of Chapter~\ref{ch:giant}; the whole of
Part~III --- the Potts correspondence, the general scheme, Ising, the hitting
set problem, vertex cover, colouring and satisfiability --- and the whole of
Part~IV, where overlapping complexes are measured and repaired. Where a chapter or a section
is new, it says so at its opening rather than leaving the reader to infer it,
and where a result of the two earlier papers is corrected here, the correction
is stated as one.
What is new in them is the formulation and not every number they return: where a chygraph calculation reproduces a published
threshold for a graph or a flat hypergraph, the agreement is a check on the
implementation and the chapter says so in those words. The thresholds that
Chapters~\ref{ch:colouring} and~\ref{ch:satisfiability} benchmark against are
other people's, attributed at each use, and the last sections of those chapters
compute them rather than quoting them. That is the check behind the one place
where the same machinery is turned on a clustered network and returns a number
nobody has computed before. All of the results can be reproduced from the
software package below. Every number the book computes comes out of it,
and the handful quoted from published work for comparison is attributed where it
appears; \emph{The software}, at the back of the book, says which routine
computes which equation, and which test checks it.

\medskip
\noindent\textbf{Key new results:}

\begin{itemize}\setlength{\itemsep}{2pt}
\item Full generating-function formalism for percolation on chygraphs,
  Chapter~\ref{ch:giant}.
\item Epidemic spreading on higher-order structures,
  Chapter~\ref{ch:epidemics}.
\item General cavity recursion for models on a chygraph,
  Chapter~\ref{ch:statmech}.
\item One-step replica symmetry breaking on a chygraph,
  Sec.~\ref{sec:onestep}.
\item Ising model on graphs with triangles, Chapter~\ref{ch:ising}.
\item Vertex covering on chygraphs, Chapters~\ref{ch:hittingset}
  and~\ref{ch:cover}.
\item Colourability of graphs with triangles, Chapter~\ref{ch:colouring}.
\item Satisfiability with nested clauses, Chapter~\ref{ch:satisfiability}.
\item Chygraphs with meta-complexes, Chapter~\ref{ch:metacomplex}.
\item \url{https://github.com/av2atgh/chygraph}.
\end{itemize}

\vspace{1em}
\noindent Alexei Vazquez
